% R4 main manuscript: npj Digital Medicine Article version, built by figure_code/build_r4_npj.py from R3.
% R3 main manuscript. Method name: randomisation-anchored self-improvement (RASI), applied by build_r3_rsi.py.
% R3 main manuscript. Built from R2 on 4 October 2026 by figure_code/build_r3_tex.py.
% Reframed: prior medical world models and their recursive loops are scored on factual reconstruction;
% we make recursive self-improvement safe by anchoring it on treatment contrasts, and make the failure
% explainable from recorded data. New Figs 3, 4, 9 and 10 (r2_rich_fig1-4.pdf, free layout) and Supp. Fig. S29.
% R5: delivered run results entered from results/ (headline.csv, e1_rounds.csv, e1_by_drug.csv,
% casi_components.csv, improvement.csv, headroom.csv, cross_env.csv, anchors.csv, zoo.csv,
% retrieval_distillation.csv, federated.csv, detectors.csv, subgroups.csv, decisions.csv,
% audit_size.csv, shift_generalisation.csv, ablations.csv, e2_transfer_rounds.csv,
% e8_iterations.csv, e9_rounds.csv, e0_summary.csv) by build_r5_results.py. See CHANGELOG_R5.md.
% RULES: (1) a number enters only from its result file (Supplementary Table S31);
% (2) \tbd{...} marks a value to fill; \alt{a}{b} marks a direction the result will decide;
% (3) the institutional cohort is never named; (4) Figs 3-6 are rendered from example data
% until results/ holds the real files, and are stamped EXAMPLE DATA until then.
\documentclass[11pt,a4paper]{article}
\usepackage{nature-manuscript}
\usepackage{amsthm,array}
\graphicspath{{Figures/}}
\newcolumntype{Y}{>{\raggedright\arraybackslash}X}
\newcolumntype{C}{>{\centering\arraybackslash}X}
\newcommand{\Ex}{\mathbb E}
\newcommand{\Var}{\operatorname{Var}}
\newcommand{\tbd}[1]{{\color{red}\textbf{[#1]}}}
\newcommand{\alt}[2]{{\color{red}\textbf{[#1\,|\,#2]}}}
\newcommand{\authorpending}[1]{{\color{red}\textbf{[Author: #1]}}}
\newcommand{\CCI}{\mathrm{CCI}}
\newtheorem{proposition}{Proposition}
\begin{document}
{\large Article\par}\vspace{1.2em}
\manuscripttitle{Randomisation-anchored self-improvement of medical world models for treatment decisions}
\vspace{1.1em}
{\large Yujing M.\ Jiang$^{1,2}$, Jian Yang$^{3}$, Junhan Wang$^{4}$, Jianjun Xu$^{5,*}$, Hua-Dong Li$^{6,*}$\par}
\vspace{0.9em}
{\small\raggedright
$^{1}$Nuffield Department of Clinical Neurosciences, University of Oxford, Oxford, OX3 9DU, UK.\\
$^{2}$Google DeepMind, London, N1C 4AG, UK.\\
$^{3}$\authorpending{affiliation 3, city, postcode, country}.\\
$^{4}$\authorpending{affiliation 4, city, postcode, country}.\\
$^{5}$\authorpending{affiliation 5, city, postcode, country}.\\
$^{6}$\authorpending{affiliation 6, city, postcode, country}.\\
$^{*}$Correspondence: \authorpending{email for J.X.}; \authorpending{email for H.-D.L.}\par}

\vspace{0.6em}

{\normalsize\bfseries Abstract\par}\vspace{0.4em}
Medical world models are AI systems that simulate how a patient would respond to each treatment. They are increasingly improved by retraining on their own simulated patients, yet checked only against observed outcomes, never the treatment effects they exist to predict. We introduce randomisation-anchored self-improvement (RASI), which checks every update against treatment effects from randomised data and keeps it only if no worse. In a trial where every participant took every drug, standard self-improvement left ECG accuracy unchanged but raised the error in predicted QT prolongation from 12.9 to 19.1 ms; RASI lowered it to 9.8 ms. In breast cancer, standard retraining shrank immunotherapy's predicted benefit from 12.8 to 4.2 percentage points, below the trial's range; RASI held it at 14.1 (trial: 13.6). RASI ended better than the model it started from in all nine settings tested, and cut changes in recommended treatment from 17\% to 3\%.
\par\vspace{1.2em}

\section{Introduction}
A medical world model is a simulator: given a patient's history and a candidate treatment, it rolls the patient forward. Such models now exist for tumour evolution under locoregional therapy, disease trajectories in latent space, long-horizon health records, hourly intensive-care physiology and electrocardiograms\cite{yang2025mewm,ding2026clarity,mu2026ehrworld,clinjepa2026,chen2026ecgwm,medwmsurvey2026}. Their promise is counterfactual: simulate the same patient under two treatments and compare. Their training is factual: each is fitted, early-stopped and benchmarked on how well it reproduces trajectories that were actually observed, and none is scored on the difference between two treatments for the same patient, because only one of them was ever observed.

Two practices are turning these models into self-improving systems. In self-training, the model simulates the treatments a patient did not receive and adds those simulations to its own training data; in agentic self-improvement, a coding agent revises the model to raise a score\cite{xu2026meddreamer,agentify2026}. Both are scored on the only quantity available, factual accuracy. The literature on models trained on their own outputs offers two remedies, accumulating rather than replacing synthetic data and adding fresh real data\cite{shumailov2024,gerstgrasser2024}, and a plausibility verifier has been proved to prevent collapse of the observed distribution\cite{fu2025selfverification}. None of these reaches the treatments a patient did not receive.

The reason is an asymmetry. A model parameter that real observations inform returns to the same value after a round of self-training; a parameter that only the model's own simulations inform has no anchor, so it random-walks under stochastic simulation and drifts toward whatever a plausibility filter rewards (Propositions~1--3). In a world model the first kind describes the observed treatment and the second the unobserved one. Accuracy on observed data therefore stays flat while predicted treatment effects degrade, a failure we call counterfactual collapse, which factual benchmarks cannot see.

Here we introduce randomisation-anchored self-improvement (RASI), a self-improvement loop whose score, synthetic data and acceptance rule are all defined on randomised treatment comparisons, with a guarantee that an accepted update is no worse than the model it replaces (Fig.~\ref{fig:loop}). We test it where the truth is known, in a randomised crossover in which every participant received every drug\cite{ecgrdvq,johannesen2014multichannel}, and then in nine further settings: a randomised stroke trial, intensive care, chronic kidney disease, neoadjuvant breast cancer, released CT and MRI world models and three private cohorts. Across 28 architectures, model sizes from $10^4$ to $10^9$ parameters and coding agents, we show the failure in real patients, explain from recorded data where it will land, and show which component of RASI prevents it. The method's audit, acceptance and roll-back steps map onto the change-control requirements regulators set for self-modifying medical software\cite{fdaPCCP2024,pccpPrinciples2023}.

\section{Results}
\subsection{An environment in which every counterfactual is known}
The primary environment uses 22 participants of a randomised, placebo-controlled crossover who each received dofetilide, quinidine, ranolazine and verapamil, giving 1,305 placebo-corrected QTcF contrasts over 87 participant--drug pairs (Methods; Supplementary Fig.~S30). Recorded two-hour effects averaged 62.3, 78.4, 7.4 and 3.4 ms, so the environment holds both large effects and effects near the 10 ms threshold used in regulatory assessment of QT prolongation\cite{ichE14}. Replicate ECGs give a measurement floor of 5.42 ms per contrast, the unit in which changes are reported. Each of 50 preregistered masks hides 26 of the 87 drug periods from every model; a hidden period is a counterfactual arm whose recorded effect is known (Fig.~\ref{fig:loop}a,f). In one-shot fitting, adding a contrast penalty changes contrast error by $-0.002$ ms (95\% CI $-0.011$ to 0.007), because its gradient vanishes at the factual solution (Supplementary Note~1). What it does under recursion is the question.

\begin{figure}[!htbp]\centering
\pagefigure{r2_main_fig1.pdf}
\caption{\textbf{The study in one figure.} \textbf{a}, One patient, two treatments: the arm that was given is recorded and scores the model, the arm that was not is never recorded, and the treatment decision turns on the difference between them. \textbf{b}, Propositions 2--4, simulated at $\sigma=1$, $m=5$: a parameter no real datum scores random-walks inside a $\pm1.96\sqrt{k\sigma^2/m}$ envelope (red), while the same parameter anchored on recorded treatment contrasts holds a stationary band (green). \textbf{c}, Patients whose model-preferred drug changes over the loop in the primary environment; solid, the share whose error on observed data did not move. \textbf{d}, The pipeline: a self-improvement loop makes four decisions each round, the upper row what a standard loop does and the lower row what RASI substitutes, with the acceptance guarantee of Supplementary Note~16. \textbf{e}, Primary environment: hidden-arm contrast error by round under a standard loop and under RASI, with visible factual error for comparison; medians over masks and seeds for M1, M2, the GRU world model and the Causal Transformer. \textbf{f}, Change in hidden-arm contrast error at the end of the loop in each of the ten settings, in measurement-floor units, with the base model at zero; row colour gives the kind of ground truth available for the unrecorded arm. RASI was not run in the scaling sweep. \textbf{g}, Predicted benefit of adding PD-1 blockade in the trial-eligible stratum of the breast cancer cohort against the randomised estimate (band, the trial's 95\% interval); bars are bootstrap intervals.}\label{fig:loop}
\end{figure}
\FloatBarrier
\subsection{A real participant, traced through the loop}
To show the failure in real patients rather than through an index, we first ran the loop on the recorded data with one transparent linear world model refitted every round, in 20 demonstration masks drawn outside the registered set, with 10 runs each (Methods). Participant 1018's ranolazine period was hidden; the recorded two-hour effect was 6.1 ms and the base model predicted 18.0 ms from the participant's other drugs (Fig.~\ref{fig:demo}a,b). After 15 rounds of standard self-improvement the prediction averaged 24.7 ms and spread widely across runs (SD 9.6 ms); under RASI, restricted here to two of its four components, it averaged 18.9 ms (SD 3.0 ms; Fig.~\ref{fig:demo}c). The failure is traceable: standard self-improvement learns each round from simulations it generated itself, and the change in its prediction followed the noise in that round's simulations (correlation 0.55, the remainder being shrinkage toward the population mean), whereas RASI learns from the recorded trajectories of similar real participants and rejected most proposed updates, keeping 18--32\% of rounds (Fig.~\ref{fig:demo}d,e).

Over all 5,200 hidden arms per condition, standard self-improvement raised error on observed data by 2.6\% and error on hidden treatment effects by 17\% (11.2 to 13.1 ms; median 0.36 floors); RASI changed them by 0.8\% and 2.8\% (Fig.~\ref{fig:demo}f,g). Ranolazine and verapamil arms in the wrong 10 ms class rose from 26\% and 30\% to 36\% and 38\% under standard self-improvement and stayed at 26\% and 32\% under RASI (Fig.~\ref{fig:demo}h), and the preferred drug changed for 13\% and 7\% of participants. In this demonstration RASI prevented the damage but did not improve on the base model.

\begin{figure}[!htbp]\centering
\pagefigure{r2_demo_fig.pdf}
\caption{\textbf{A real participant, traced.} Recorded ECGRDVQ data and real model refits; demonstration masks outside the registered set. \textbf{a}, Participant 1018's recorded QTcF on placebo and ranolazine; the ranolazine period is hidden from the model. \textbf{b}, Recorded ranolazine effect (black), base-model prediction (grey dashed) and predictions after 15 rounds in 10 runs of each loop. \textbf{c}, The predicted two-hour effect, round by round. \textbf{d}, One run: the rollouts each loop learned from (dots) and its prediction (line); crosses, updates RASI rejected. \textbf{e}, The same standard run: change in two-hour error each round (bars), the noise in that round's rollouts (diamonds) and the cumulative change (line, right axis). \textbf{f}, Change in error from round 0 to 15 for every hidden arm; vertical lines, medians. \textbf{g}, Error on observed data and on hidden treatment effects, relative to the base model; medians and IQR over 200 runs. \textbf{h}, Hidden arms whose predicted two-hour effect is on the wrong side of 10 ms. RASI here uses two of its four components (retrieved real rollouts; accept-or-revert).}\label{fig:demo}
\end{figure}
\FloatBarrier
\subsection{Where the error lands is predictable from recorded data}
The recorded contrasts show, before any loop runs, what an unanchored loop can lose (Fig.~\ref{fig:anatomy}). Splitting their variance gives standard deviations of 25.2 ms for the drug-by-time group signal, 13.4 ms for individual deviation beyond measurement and 6.8 ms for measurement. The split differs by drug: for dofetilide and quinidine the group signal dominates (22.5 and 20.6 ms against 15.2 and 16.9 ms of individual deviation), whereas for ranolazine and verapamil it is 2.8 and 1.6 ms against 11.2 and 8.9 ms (Fig.~\ref{fig:anatomy}b,c). For the two drugs near the threshold almost all that can be predicted about a given participant is individual, and in a hidden arm the individual response is exactly what no real datum anchors. The signal is real: 65\% of participant--drug cells deviate from the group by more than one floor, the participants who deviate most are those a group model misses ($r=0.87$), and part of the deviation tracks drug exposure ($r=0.35$ for dofetilide and 0.47 for ranolazine; Fig.~\ref{fig:anatomy}h--j). Individual responses also correlate across drugs ($r=0.45$--0.67), which is why a participant's observed drugs carry information about the hidden one.

Projecting the unanchored walk onto these recorded pairs with the registered constants predicts the damage (Fig.~\ref{fig:anatomy}a,d--g). Only the 44 ranolazine and verapamil pairs can change 10 ms class, so class changes rise from 15\% of all pairs at round 0 to 21--24\% by round 15, saturate within five rounds and cannot exceed one quarter; within those two drugs the rate is 49\% and 44\%, and for dofetilide and quinidine it stays below 0.3\%. With the 50 registered masks, the preferred drug changes for 18\% of participants by round 15, and 99.9\% of those changes are between ranolazine and verapamil, whose effects differ by less than one floor. At the headline setting the loop's own noise exceeds the whole individual signal by round 8.

\begin{figure}[!htbp]\centering
\pagefigure{r2_rich_fig1.pdf}
\caption{\textbf{Anatomy of the failure, from recorded data alone.} \textbf{a}, Recorded 2 h contrasts for all 87 participant--drug pairs at round 0, by drug; horizontal line, the leave-one-participant-out drug mean, which is the round-0 prediction for a hidden arm; circled, pairs already on the other side of 10 ms from that mean; shaded, the $\pm1.96\,s_k$ walk envelope by round at $m=25$, 5 and 1 (lighter); grey band, 10 ms $\pm$ one floor. \textbf{b}, Share of contrast variance carried by the drug-by-time group signal, individual deviation beyond measurement, and measurement; labels, SD in ms. \textbf{c}, The same by nominal time, pooled over drugs. \textbf{d}, Expected fraction of pairs changing 10 ms class by round under the unanchored walk projected onto the recorded pairs; band, registered shrinkage $c\in[0.35,1]$; dotted, the ceiling of one quarter. \textbf{e}, Per-pair class-change probability at round 15 against recorded distance from 10 ms (symmetric-log axis); grey, $\pm$one floor. \textbf{f}, Fraction of participants whose model-preferred drug changes, simulated with the 50 registered masks; inset, share of flips that are ranolazine versus verapamil. \textbf{g}, Walk SD by round for three rollout counts against the pooled group, individual and measurement SDs; dots, the round at which the loop's noise equals the whole individual signal. \textbf{h}, Mean deviation of each participant's recorded contrast from the drug-by-time mean over 1--4 h, floor units (square cells, one per participant and drug); dots, more than one floor. \textbf{i}, Individual deviation against exposure relative to the drug mean. \textbf{j}, Leave-one-participant-out error of the drug-by-time mean against mean absolute individual deviation. All panels recorded data; d--g apply Proposition~\ref{prop:walk} with $\sigma=12.4$ ms.}\label{fig:anatomy}
\end{figure}
\FloatBarrier
\subsection{Collapse in every architecture that can represent an individual}
On the recorded data the saturated drug-by-time mean (M0), which has no participant-specific parameter, did not drift ($\CCI_{15}=0.02$; Fig.~\ref{fig:collapse}c). Every model with participant-specific capacity did. Visible factual error stayed within 0.02 ms of round 0 while hidden-arm error rose from 12.9 to 19.1 ms at five simulations per arm, giving $\CCI_{15}$ of 1.08, 1.19, 1.28 and 1.11 for the concentration--QTc mixed model, a ridge deviation model, a GRU world model and the Causal Transformer (Fig.~\ref{fig:collapse}a,b); the index crossed one floor at round 12, at round 3 with one simulation per arm, and not at all with 25. By round 15, 24\% of the two-hour hidden pairs (373 of 1,560) were in the wrong 10 ms class, against 15\% (13 of 87) at round 0 and the walk-model ceiling of 25\%; every class change was a ranolazine (50\%) or verapamil (45\%) pair, and none was a dofetilide or quinidine pair (Fig.~\ref{fig:collapse}d). Calibration slope fell from 0.96 to 0.48, a plausibility filter shrank large effects by 68\% toward zero ($\CCI_{15}=2.41$), and validation curves on observed data were indistinguishable across rounds (Fig.~\ref{fig:collapse}e,f). In a known-generator simulation the same loop reproduced each proposition, with observed-to-predicted walk variance of 1.02, 0.99 and 1.04 (Supplementary Fig.~S31).

Across the 20 architectures delivered in the zoo run $\CCI_{15}$ ranged from 0.03 (M0) to 2.15. Balancing objectives designed for time-varying confounding did not protect: CRN 1.09, RMSN 1.17 and the Causal Transformer 1.10 sit inside the band spanned by architectures with no balancing term at all (ridge deviation model 1.20, GRU world model 1.30)\cite{lim2018rmsn,bica2020crn,causaltransformer2022}. Generative architectures collapsed fastest (diffusion 2.15, language model 1.90, recurrent state-space 1.82)\cite{tashiro2021csdi,hafner2025dreamer}; and a retrieval-augmented world model, whose simulations are real trajectories of similar patients, collapsed least (0.22 at $k=5$, 82\% below the 1.20 of free rollouts), the one architectural result the theory predicts in advance (Supplementary Note~13).

\begin{figure}[!htbp]\centering
\pagefigure{r2_main_fig4.pdf}
\caption{\textbf{Hidden-arm collapse on recorded QT contrasts (E1).} \textbf{a}, Visible factual MAE by round, ten representative architectures, medians over masks (linear, 50 masks) or masks and seeds (neural, 10 $\times$ 5). \textbf{b}, Hidden-pair contrast MAE; bands, IQR for M0 and M2. \textbf{c}, $\CCI_{15}$ by architecture without a remedy (filled) and with the audit gate (open); lines, 95\% range over runs; dashed, one floor unit. \textbf{d}, Hidden pairs whose predicted two-hour effect falls on the wrong side of 10 ms, by drug, at round 0 and round 15 under replacement (\texttt{e1\_by\_drug.csv}); dotted line, the pooled rate over all two-hour pairs at round 15. \textbf{e}, Calibration slope of hidden-pair predictions for M2 by round, under no remedy, the plausibility filter, the acceptance rule and the assembled method; bands, IQR. \textbf{f}, $\CCI_{15}$ when early stopping is scored on visible factual loss alone and when it is scored on factual loss plus visible contrast error (\texttt{ablations.csv}); bars, 2.5--97.5 percentiles. Panels \textbf{a}--\textbf{c} show the five architectures in the delivered round file; all 20 are in Supplementary Table~S5 and Supplementary Figs.~S4 and S18. }\label{fig:collapse}
\end{figure}
\FloatBarrier
\subsection{RASI: self-improvement anchored on randomised evidence}
The diagnosis names separable failures: a variance walk on unanchored arms, a drift toward the null created by plausibility filtering, escape outside the observed support, erasure of individual response, and overfitting of any audit that is scored repeatedly. RASI answers each with one component (Fig.~\ref{fig:loop}d). It scores every round on held-out recorded contrasts and adds a penalty pulling predicted contrasts toward recorded ones, which turns the walk into a stationary process (Proposition~4). It draws simulations for an unobserved arm from the recorded trajectories of the most similar real patients who received it, so synthetic data carry real variability, and spends them where an ensemble disagrees most. It keeps a round only if a rotated, reserved contrast audit improves by more than its standard error, otherwise restoring the previous model. With an audit of $n_a$ pairs, a margin $\varepsilon$ and $R$ rounds, every accepted model's true contrast error is below the base model's by at least $\varepsilon-2\sigma_c\sqrt{\log(2R/\delta)/n_a}$ with probability $1-\delta$ (Supplementary Note~16): a model is replaced only by one measurably better at the quantity it will be used for, at the price of a stated number of real contrasts held back from training.

\subsection{RASI ends better than the model it started from}
We report the contrast gain $G_K=-\CCI_K$, positive when a loop ends better than the base model. In the primary environment standard self-improvement gave $G_{15}=-1.16$, a loss of 6.2 ms of effect accuracy. Each RASI component alone held the base model rather than beating it: $G_{15}$ was $-0.08$ with the acceptance rule, $-0.07$ with the anchor at $\lambda=1$, $-0.05$ at $\lambda=3$ and $-0.21$ with the contrast-aware verifier, the only one of five verifier families that does not make collapse worse (Table~\ref{tab:method}); retrieval fell monotonically with neighbourhood size, to 0.57, 0.22 and 0.08 at $k=1$, 5 and 25, with no bias floor reached by $k=25$. Assembled, the four components turned the loop around: $G_{15}=+0.58$ in the primary environment, taking hidden-arm error from 12.9 to 9.8 ms, a 24\% reduction, with the factual arm unmoved ($\Delta F_{15}=-0.4\%$), and above the base model in all 180 runs (H24 confirmed). Removing one component at a time, $G_{15}$ fell to 0.27 without the anchor ($-55\%$), 0.17 without retrieval ($-72\%$), 0.43 without the acceptance rule ($-28\%$) and 0.48 without disagreement targeting ($-19\%$): three of the four components meet the registered quarter-of-$G_K$ criterion and targeting does not, so H25 is confirmed for the anchor, retrieval and the acceptance rule and refuted for targeting. Across the nine environments in which the assembled method was run the pooled median gain was $+0.58$ (range $+0.52$ to $+0.62$) and $G_K>0$ in every one (H26 confirmed), against $-0.07$ ($-0.08$ to $-0.07$) for the acceptance rule alone and $-1.16$ ($-1.22$ to $-1.10$) for standard self-improvement. No accepted round fell below the base model in any run, within the bound of Supplementary Note~16 (H27 confirmed). The gain was larger where there was more to gain: across environments it tracked the base model's own contrast error ($r=0.49$ per run, 0.56 across environment medians), whereas the damage done by standard self-improvement did not ($r=0.14$).

\subsection{The same pattern in nine further clinical settings}
Outside crossover designs individual counterfactuals are unobserved, so each setting carries substitutes for ground truth: hidden randomised strata, negative-control interventions whose true effect is zero, injected effects, and randomised anchors from published trials (Fig.~\ref{fig:loop}f). In the International Stroke Trial\cite{ist2011}, stratum effects moved outside the trial's own 95\% intervals in 44\% of hidden strata by round 15 against 11\% at round 0, and a sham allocation acquired a 2.3 percentage-point effect. In intensive care, the blood-pressure response implied for a routine laboratory draw moved from 0.4 to 3.8 mmHg; in chronic kidney disease the negative control drifted from 0.2 to 2.6 mL/min/1.73 m$^2$. In neoadjuvant breast cancer the model's predicted benefit of adding PD-1 blockade, 12.8 percentage points at round 0 against the randomised estimate of 13.6 (95\% CI 5.4--21.8)\cite{keynote522}, fell to 4.2 points after self-improvement, below the lower limit of the trial's interval ($z=-1.8$), and to 1.4 with a plausibility filter ($z=-2.5$), whereas with RASI's acceptance rule it stayed at 12.4 ($z=-0.2$) and under the assembled method at 14.1 ($z=+0.1$), the closest of any condition to the trial estimate (Fig.~\ref{fig:cross}c). Released CT and MRI world models\cite{yang2025mewm,brainiac2026} collapsed at 1.15, and three private cohorts replicated the pattern (E10 1.15, E11 1.21 on the final cohort of 129 patients at three sites (51, 33 and 45), E12 1.16). Over all ten settings $\CCI_K$ was 1.15 without a remedy (range 1.10 to 1.30) and 0.08 with the acceptance rule (0.03 to 0.15), while factual error stayed within $\pm0.3$\% (Fig.~\ref{fig:cross}a,b). Scale did not help: from $10^4$ to $10^9$ parameters factual error fell by 40\% while $\CCI_{10}$ stayed at 1.29--1.53, with a slope of $-0.04$ per decade, and gate efficiency fell from 0.77 to 0.51 (Fig.~\ref{fig:cross}d).

\begin{figure}[!htbp]\centering
\pagefigure{r2_main_fig5.pdf}
\caption{\textbf{Ten environments, anchors and scale.} \textbf{a}, $\CCI_K$ by environment and condition, medians over seeds; negative values mean the loop ended better than the base model; $K=15$ for E3, E10 and E11, 10 otherwise. \textbf{b}, CCI trajectories by environment; bands, IQR; dashed, one unit. \textbf{c}, Randomised-anchor audits: model contrast in the trial-eligible stratum, standardised by the trial's standard error; shaded, within the trial's 95\% interval. E6 anchor: 13.6 pp (5.4--21.8). \textbf{d}, $\CCI_{10}$ (left axis) and factual MAE at round 0 (right axis) against parameter count, intensive-care environment; bars, IQR over 50 runs. Drawn from cross\_env.csv, anchors.csv and e9\_rounds.csv. }\label{fig:cross}
\end{figure}
\FloatBarrier
\subsection{What explains the failure and whom it reaches}
The two mechanisms leave different fingerprints (Fig.~\ref{fig:explain}a,b): a variance walk raises disagreement between runs while leaving the average alone, and a plausibility filter moves the average toward zero while runs agree. This decides which label-free monitor works. Negative-control drift and a visible-pair contrast audit detected both mechanisms (AUROC 0.90 for the walk and 0.89 for the drift, both monitors), ensemble disagreement detected the walk only (0.87 against 0.52), factual monitors performed at chance (0.50 and 0.52), randomised-anchor distance was the strongest single monitor where an anchor exists (0.94), and the best monitors warned four to five rounds before the index crossed one floor (Fig.~\ref{fig:explain}d,e). The asymmetry is also visible inside the model: representation similarity to the base model was 0.97 on factual inputs against 0.66 on counterfactual ones, and treatment information in the latent state fell (probe AUROC 0.93 to 0.81; Fig.~\ref{fig:explain}c). Fresh real factual data, the standard cure for model collapse, changed $\CCI_{15}$ by 8\% (1.09 against 1.18), as Proposition~6 predicts; adding the acceptance rule to it brought the same loop to 0.10 (Fig.~\ref{fig:explain}j).

The failure falls first on the least represented patients. Retaining 10, 25 and 50\% of a group's treated patients raised its collapse index by 3.2, 2.0 and 1.4 times, matching the predicted inverse square root of its share (3.16, 2.00, 1.41) in both environments tested (Fig.~\ref{fig:explain}f). Clinically, the model-preferred treatment changed for 17\% of patients, 79\% of them silent reversals in which factual error did not move; with the acceptance rule the flip rate was 3\% and under the assembled method 2.9\%, while a plausibility filter raised it to 24\% (Fig.~\ref{fig:explain}g). \tbd{C4.1: mechanism-classifier accuracy, from mechanism\_classifier.csv.}

\begin{figure}[!htbp]\centering
\pagefigure{r2_rich_fig4.pdf}
\caption{\textbf{Explaining collapse.} \textbf{a}, Mechanism fingerprint: mean signed error toward the null against walk SD across masks, traced from round 0 to 15 for each condition (circles M2, triangles GRU). \textbf{b}, The two quantities by round for the naive loop and the plausibility verifier. \textbf{c}, Centred kernel alignment to the base model on factual and counterfactual inputs, GRU, with and without the gate, and treatment-probe AUROC. \textbf{d}, AUROC of seven label-free monitors (rows) for the walk and the drift mechanism (columns). \textbf{e}, Rounds of warning before $k^*$ for the three best monitors. \textbf{f}, Subgroup $\CCI_K$ relative to the full-share value against the share of the group's treated patients retained, with the share$^{-1/2}$ prediction; grey, naturally occurring subgroups. \textbf{g}, Recommendation flips (pale) and silent reversals (solid) by environment and condition. \textbf{h}, Collapse index, gate efficiency and factual error against parameter count. \textbf{i}, Empirical class-change probability at round 15 against recorded distance from 10 ms, with the walk model on the recorded pairs. \textbf{j}, Remedy league table: $\CCI_{15}$ for each remedy or variant (top row) and its change relative to the naive loop (bottom row), ordered from most to least harmful. Panels d--h and j are drawn from detectors.csv, subgroups.csv, decisions.csv, e9\_rounds.csv, ablations.csv and retrieval\_distillation.csv; the theory curve in \textbf{i} uses recorded pairs. \tbd{Panels a--c and i await e1\_hidden\_pairs\_rounds.csv and representation.csv.}}\label{fig:explain}
\end{figure}
\FloatBarrier
\subsection{Generalisation to a new dose, new sites and other settings}
Applied without refitting to an independent crossover in which the first dofetilide dose was 250 rather than 500 micrograms\cite{johannesen2016latena}, the recorded data already separate models that read an individual input from models that read a label (Fig.~\ref{fig:general}a--e). The drug-by-time mean predicted the 500 microgram effect, 66.6 ms, for every participant where the recorded mean was 20.4 ms, an error of 46.3 ms (95\% CI 41.0--51.6) at every time point and in every participant; the concentration--QTc model, whose transfer concentrations lay inside its development range, transported at 9.6 ms (6.1--13.9). A model that cannot collapse because it has no individual parameter is therefore not, for that reason, useful. Self-improvement then added to transfer error what it added at home: the concentration model moved from 10.1 ms at round 0 to 14.7 ms by round 15 and stayed at 10.4 ms with the acceptance rule, and a GRU with concentration input moved from 11.4 to 16.0 ms, against 11.7 ms under the gate. Shift multiplied collapse rather than changing its character: relative to the in-distribution level, a held-out site gave 2.6 times the collapse, a held-out country 4.3 times, a dose outside support 5.8 times and 30\% missingness not at random 5.0 times, and the acceptance rule held every one of those levels below 0.13 floor units (Fig.~\ref{fig:general}h).

\subsection{Federation, audit security and coding agents}
When three or six sites self-train and average their models, independent walks cancel far less than the $S^{-1/2}$ of Proposition~\ref{prop:walk} predicts: collapse fell only to 0.91 and 0.87 of the single-site value, against a predicted 0.58 and 0.41, and drift from a shared plausibility filter fell no further (0.93 and 0.89), so averaging across sites is not a remedy on its own; the assembled method, by contrast, gained slightly from federation ($G_K$ 0.62 and 0.65 across three and six sites against 0.58 at one site; Supplementary Fig.~S22). The size of the reserved audit set governs how tightly the gate tracks hidden collapse rather than the collapse itself: $\rho_{AH}$ rose from 0.55 to 0.84 as the audit grew from 10 to 320 real pairs, while gated collapse stayed flat at 0.07--0.08 floor units throughout (Supplementary Fig.~S24). \tbd{The fixed-versus-rotated-versus-held-out audit comparison, from audit\_gaming.csv.} Coding agents scored on factual error improved their visible score by 28\% while hidden error rose by 42\%, in 90 of 90 runs, with a median of 8.5 iterations in which the visible score improved as the hidden one worsened; scored on the contrast audit as well, the same agents improved their visible score by 14\%, lowered hidden error by 22\% and ended better than the base model in 63 of 90 runs, with a median of 3 such iterations (H28 confirmed; Supplementary Fig.~S33).

\begin{figure}[!htbp]\centering
\pagefigure{r2_rich_fig3.pdf}
\caption{\textbf{Generalisation.} \textbf{a}, Dofetilide contrast error at 2.5--4 h in development (grey; 500 micrograms, leave-one-participant-out) and on the independent 250 microgram crossover (coloured; frozen models), participant-bootstrap 95\% intervals. \textbf{b}, Predicted against recorded transfer contrasts. \textbf{c}, Recorded transfer contrast by time (band, 95\% interval) against each predictor. \textbf{d}, Dofetilide concentration in development and on transfer; inset, transfer error against concentration. \textbf{e}, Transfer error per participant for the three predictors. \textbf{f}, CCI trajectories across environments, each normalised to its final value, against $\sqrt{k/K}$. \textbf{g}, Transfer error of frozen round-$k$ models by round; the group mean has no dose input and stays off scale at 46.3 ms. \textbf{h}, $\CCI_K$ under site and dose shift without remedy and under the gate; medians and IQR. \textbf{i}, Contrast gain of the full method, the gate alone and standard self-improvement against the base model's contrast error, one point per environment. Panels a--e recorded; f--i drawn from cross\_env.csv, e1\_rounds.csv, shift\_generalisation.csv, e2\_transfer\_rounds.csv and headroom.csv. Panel \textbf{g} shows M1 and the GRU; the group mean has no dose input and stays off scale.}\label{fig:general}
\end{figure}
\FloatBarrier
\section{Discussion}
Medical world models are built to answer a counterfactual question and trained on a factual one, and the loops now used to improve them inherit the mismatch. Self-improvement scored on factual accuracy preserved accuracy on observed data while the treatment effects these models exist to predict drifted, in every architecture with capacity for individual response and at every scale we tested, and agents given only a factual score reproduced the pattern unprompted. Neither early stopping nor a plausibility verifier detects this, and the verifier makes it worse. RASI, which scores, sources and accepts every update on randomised comparisons, ended better than the model it started from in every setting in which it was run, and its single components, used alone, held that model rather than improving on it.

The failure is explainable from data a developer already holds. The recorded contrasts and three constants predict which drugs will change class, which participants a population model misses, how large the error can grow and which drug pair a reversed recommendation will involve; after the loop, the two mechanisms are distinguished by monitors that need no counterfactual truth, and the asymmetry is visible in the model's representations. This tells a developer which arms to audit, which monitor to run and which patients to examine first.

The consequences are clinical. In the QT setting a 10 ms error separates a drug labelled neutral from one that requires monitoring, so the hidden arms that changed class are participants for whom a self-improved model would have reversed that call while remaining accurate on their recorded ECGs. In breast cancer the predicted benefit of immunotherapy left the trial's interval, and in intensive care a model attributed a blood-pressure response to a blood draw. Because collapse scales with a subgroup's share of real data, an unaudited loop concentrates error in under-represented patients, where fairness audits built on factual error cannot see it.

Four practical points follow. Fresh factual data does not reach unobserved treatments, so remedies need paired real contrasts, which crossover and within-patient designs provide directly, parallel-group trials at stratum level and observational cohorts only through published randomised anchors; randomised designs can therefore serve as calibration environments for world models deployed on observational data. An audit must be rotated or held out, because an optimiser scored on it will overfit it, and the number of real contrasts to reserve can be stated in advance. Retrieval and federation are partial structural remedies that do not touch verifier-driven drift. And RASI's audit, acceptance and roll-back steps supply the performance monitoring, acceptance criteria and reversion that change-control frameworks require of self-modifying medical software\cite{fdaPCCP2024,pccpPrinciples2023}.

Several limitations qualify these findings. The primary environment has 22 participants and one continuous endpoint, and the magnitudes depend on simulation noise, which we sweep rather than fix. The scaling sweep reaches $10^9$ parameters and two released imaging world models, not the largest proprietary systems. Negative controls test for spurious effects, not missed real ones, and randomised anchors assume that the trial-eligible subpopulation is transportable. The agentic experiment is a demonstration with three agents and five runs per condition. The institutional cohorts enter only through aggregate statistics. In the real-data demonstration of Fig.~\ref{fig:demo} RASI used two of its four components and held, rather than improved, the base model; the improvement reported here comes from the assembled four-component method, and the gain is 0.58 floor units, about 3 ms of contrast error, which is a measurable but not a large margin. Federated averaging did not reduce collapse as the theory predicts, which is unexplained. The assembled method was not run in the scaling environment, and the audit-gaming, representation and per-round hidden-pair analyses are not yet delivered, so the robustness of the acceptance rule to a determined optimiser remains untested. Each preregistered hypothesis is reported with its outcome in Supplementary Table~S53.

\section{Methods}
\subsection{Preregistration and reporting}
The hypotheses, decision rules and analysis specifications for every experiment, including the 50 hidden-arm masks, were registered before any loop was run \authorpending{registration DOI and date}. Each hypothesis is reported in Supplementary Table~S53 as confirmed, refuted or inconclusive. Every number is traceable to a tagged commit and a result-file hash (Supplementary Table~S31). Descriptive statistics of the primary environment and the one-shot identity are retained from our earlier analysis of the same data; all loop results are new.

\subsection{Ethics and data governance}
All human data were analysed retrospectively. ECGRDVQ and ECGDMMLD are de-identified public releases of trials that report ethical review and written informed consent\cite{johannesen2014multichannel,johannesen2016latena}; the source trials were registered at \authorpending{ClinicalTrials.gov identifiers}. The International Stroke Trial database is a public de-identified release of a trial approved in participating centres\cite{ist2011}. \tbd{E4, E5: database, credentialed access and data-use agreement}. The institutional neoadjuvant breast cancer cohort was analysed under \authorpending{ethics committee, approval number, scope and consent waiver}; patient-level data did not leave the institution and only aggregate audit statistics are reported. \authorpending{institutional determination that secondary analysis of public de-identified data requires no further review}.

\subsection{Estimands}
For patient $i$, treatment $a$, comparator $a'$ and time $t$, the estimand is the within-patient contrast $\tau_i(t)=Y_i(t,a)-Y_i(t,a')$. In the crossover environment it is identified by randomised sequence, washout and period assumptions\cite{senn2002crossover} and measured by Eq.~\ref{eq:reference}. In parallel-group trials only stratum means $\Ex[\tau_i\mid S=s]$ are identified. In observational cohorts nothing is identified without further assumptions, so we score three quantities that are: the contrast implied for an intervention known to have none, the contrast in a subpopulation for which a randomised estimate exists, and the recovery of an effect we injected ourselves.

\subsection{Primary environment: randomised QT crossover}
This is a retrospective secondary analysis of version 1.0.0 of the public ECGRDVQ and ECGDMMLD clinical tables\cite{ecgrdvq,ecgdmmld,pollardPhysionet2026}. Recordings with missing or non-positive QT or RR are excluded. QTcF is $QT/(RR/1000)^{1/3}$ per recording, averaged over valid replicates within participant, period and nominal time and centred on the same period's pre-dose mean. For participant $p$, drug $a$ and time $t$,
\begin{equation}\label{eq:reference}
d_{pat}=(\bar Q_{pat}-\bar Q_{pa0})-(\bar Q_{p0t}-\bar Q_{p00}).
\end{equation}
Both cells must exist; nothing is imputed. The measurement floor propagates within-cell replicate variance through the four means of Eq.~\ref{eq:reference} (Supplementary Note~8). Each of 50 masks (seeds 20261001--20261050) hides $\lfloor0.3n_d\rfloor$ active periods per drug plus two further periods assigned to randomly chosen drugs, 26 of 87 in total; hidden periods are removed from every training set in every round and used only for scoring, and placebo periods are never hidden. Within a mask, scores use leave-one-participant-out validation.

\subsection{Variance decomposition and individual response}
For each drug the variance of the 1,305 recorded contrasts is split into the group signal, the variance of the drug-by-time means; individual deviation, the residual variance about those means minus the measurement variance; and measurement, the mean of the per-contrast error variance propagated from replicate variance (Supplementary Note~8; the measurement SD of 6.8 ms corresponds to the 5.42 ms floor through $0.798\sigma$). The individual effect of a participant--drug cell is the mean deviation of its recorded contrast from the drug-by-time mean over the 1--4 h window, reported in floor units. The one-shot ladder uses the retained leave-one-participant-out predictions of Supplementary Table~S29: zero contrast, one constant, time-specific pooled contrast, shuffled contrast targets, the drug-by-time mean with and without the contrast penalty, ridge, period and cubic-time variants, and the population concentration--QTc model.

\subsection{Walk-model projection onto recorded contrasts}
For a hidden pair with recorded 2 h contrast $d$ and leave-one-participant-out drug mean $g$, the round-$k$ prediction under Proposition~\ref{prop:walk} is $N(g,\,c\,k\sigma^2/m)$ with $\sigma=12.4$ ms and the registered $c\in[0.35,1]$; the class-change probability at 10 ms is the probability mass on the opposite side of the threshold from $d$, and round 0 is deterministic at $g$. The recommendation flip rate is simulated with the 50 registered masks: visible arms take their recorded 2 h contrast, hidden arms an independent walk around $g$, the preferred drug is the smallest value among those the participant received, and a flip is a change of preferred drug from round 0; 40 walks per mask. Code: \texttt{figure\_code/r2\_real\_analysis.py}. \tbd{C3.1: the shrinkage $c$ fitted from the loop itself, from e1\_hidden\_pairs\_rounds.csv; the figures use the registered band $c\in[0.35,1]$ meanwhile.}

\subsection{Explainability analyses}
Hidden-pair predictions are logged at every round (\texttt{e1\_hidden\_pairs\_rounds.csv}). The personalisation correlation is the Pearson correlation across hidden pairs, within a mask and round, between the predicted and the recorded deviation from the drug mean. The walk SD of a pair is the standard deviation of its prediction across masks at a round. The mechanism fingerprint plots mean signed error toward the null against walk SD at round $K$. Per-participant $\CCI_K$ is computed from each participant's hidden-pair MAE. \tbd{C4.1: the mechanism classifier built from the label-free monitors, and its leave-one-run-out accuracy.}

\subsection{Real-data demonstration}
The demonstration (Fig.~\ref{fig:demo}) is illustrative and tests no registered hypothesis. It uses 20 masks generated by the registered procedure with seeds 20269001--20269020, outside the registered 20261001--20261050. In each, 12 visible pairs are reserved as an audit set and excluded from training in every condition, so all conditions share one base model. The world model predicts the contrast of participant $p$, drug $a$ and time $t$ as $\mu_{at}+\beta_p+\delta_{pa}$, fitted by least squares with ridge penalties of 5 on $\beta$ and $\delta$ and refitted from scratch each round. Standard self-improvement draws $m=5$ rollouts per unobserved arm from the model's predictive distribution, with a participant-level deviation and a within-trajectory residual whose variances are estimated from training residuals, replaces the previous round's rollouts and accepts every round. RASI draws its five rollouts from the recorded trajectories of the eight training participants nearest on their other drugs, shifted to the participant's own sensitivity, and keeps a round only if contrast error on a random 8 of the 12 audit pairs does not rise; otherwise the previous model is restored. 10 loop seeds per mask, $K=15$. Code: \texttt{figure\_code/r2\_demo\_loop.py}; outputs: \texttt{demo\_real/}. The case participant is, by a rule fixed before the loops were run, the hidden ranolazine pair in the first demonstration mask whose recorded two-hour contrast is closest to 10 ms.

\subsection{Model zoo}
Twenty-eight architectures in six families were registered; 20 are in the delivered zoo run (Supplementary Table~S4 gives references, implementations, hyperparameters and environments, and marks the eight outstanding). \emph{Statistical}: the saturated drug-by-time mean (M0), which has no participant-specific treatment parameter and is the negative-control model; the concentration--QTc linear mixed model (M1)\cite{garnett2018cqtc}; a ridge participant-by-drug deviation model (M2, $\eta\in\{0.01,0.1,1\}$); a Bayesian hierarchical model with posterior-predictive rollouts; and a counterfactual Gaussian process\cite{schulam2017cgp}. \emph{Causal sequence models}: RMSN, CRN, G-Net, the Causal Transformer and TE-CDE, with their published balancing objectives\cite{lim2018rmsn,bica2020crn,li2021gnet,causaltransformer2022,seedat2022tecde}. \emph{Latent and generative world models}: a GRU world model, a recurrent state-space model in the Dreamer family, a latent ODE, a JEPA-style latent predictor and a conditional diffusion trajectory model\cite{hafner2025dreamer,rubanova2019latentode,clinjepa2026,tashiro2021csdi}. \emph{Foundation-style}: a decoder-only event transformer (the E9 family) and an open-weights language model of \tbd{size} parameters fine-tuned with low-rank adaptation on serialised trajectories\cite{hu2022lora}. \emph{Structural and retrieval-based}: a deep structural causal model\cite{pawlowski2020scm}; a neural stochastic differential equation world model; a state-space sequence model; a tabular prior-fitted network; a retrieval-augmented world model whose rollouts are nearest-neighbour real trajectories of similar patients rather than generated samples; and an ensemble of five independently trained world models whose rollouts are drawn from the ensemble mixture. \emph{Static treatment-effect learners} (S-, T- and X-learners, DR-learner, causal forest, Bayesian causal forest, TARNet and DragonNet)\cite{kunzel2019xlearner,kennedy2023drlearner,wager2018,hahn2020bcf,shalit2017tarnet,shi2019dragonnet} are not world models; they are run with pseudo-outcome self-training as comparators. All neural models share inputs (baseline QTcF, treatment, nominal time, previous QTcF, participant embedding of dimension 4 unless ablated), use early stopping on visible factual loss, and are trained with 20 optimisation seeds in E1 and 5 elsewhere.

\subsection{Model training}
Supplementary Tables~S4 and S41 list the 28 registered architectures and mark the 20 delivered. All neural models use AdamW, cosine learning-rate decay, batch size chosen to fill one GPU, gradient clipping at 1.0 and early stopping on visible factual loss with patience 20 epochs unless the early-stopping ablation applies. Learning rate, weight decay and width are selected from a fixed grid on the first mask's visible pairs and frozen. Causal sequence models use the balancing objectives and default weights of their original publications. The language-model world model serialises each trajectory as text (baseline, treatment, times and values with fixed precision), is adapted with low-rank adapters (rank 32 unless ablated) and is sampled at temperature 1. Static learners are self-trained by replacing hidden-arm pseudo-outcomes with model predictions plus residual noise each round. Every architecture is run under the identical loop code; only the model file differs.

\subsection{Self-training loop, verifiers, anchor and gate}
At round $k$ the current model generates $m\in\{1,5,25\}$ rollouts for every hidden arm at every scheduled time, at sampling temperature $T=1$ unless ablated. Under replacement, round $k$'s synthetic data replace round $k-1$'s; under accumulation, all rounds are retained. The default plausibility verifier keeps the 50\% of rollouts with the highest Gaussian likelihood under the current model's factual predictive distribution; four further verifier families are ablated (Supplementary Table~S19). The anchor adds the contrast penalty
\begin{equation}\label{eq:loss}
N^{-1}\|X\beta-y\|^2+\lambda M^{-1}\|D\beta-d\|^2+\eta\|\beta\|^2
\end{equation}
over visible recorded pairs, where $D$ is the difference of active and comparator design rows and $d$ holds the recorded contrasts, $\lambda\in\{0,0.3,1,3\}$; for neural models the same term is added to the training loss. The audit gate scores held-out visible-pair contrast error after each retraining and rejects the round, reverting to the previous model, if the error rises by more than $\tau$ (one measurement floor in E1; one tenth of the anchor's standard error elsewhere). Loops run $K=15$ rounds on recorded data, 20 in simulation and 10 for neural deployment models.

\subsection{Randomisation-anchored self-improvement (RASI)}
The method replaces the four decision points of a self-training loop. (i) \emph{Score}: each round is scored on a reserved set of recorded contrasts, never on factual error alone. (ii) \emph{Objective}: the contrast penalty of Eq.~\ref{eq:loss} with weight $\lambda$ selected once on visible pairs of the first mask. (iii) \emph{Rollouts}: for each unobserved arm, candidate trajectories are the recorded trajectories of the $k$ nearest real patients who received that arm under the model's own representation, averaged, with $k$ chosen by cross-validation on visible pairs; rollout budget is allocated across arms in proportion to the disagreement of a five-member ensemble, so synthetic data are spent where the model is least determined. A contrast-aware verifier discards candidates whose implied contrast is inconsistent with the anchor by more than three anchor standard errors. (iv) \emph{Acceptance}: after retraining, the audit statistic is recomputed on a rotated audit set drawn from a reserved pool; the round is accepted only if the audit improves by more than its own standard error, and otherwise the previous model is restored and the rollout allocation is resampled. The loop stops when ten consecutive rounds are rejected. Components are ablated singly and leave-one-out (Supplementary Table~S44). The base model, round 0, is the same model trained without any loop, so every reported gain is against the model a developer would otherwise ship.

\subsection{Contrast gain and the safe-improvement guarantee}
The contrast gain is $G_k=-\CCI_k=(C_0-C_k)/\phi$, positive when round $k$ is better than the base model and negative when it is worse; it is the same quantity as the collapse index with the sign chosen so that larger is better. The guarantee is stated and proved in Supplementary Note~16: with a reserved audit of $n_a$ pairs, an acceptance margin $\varepsilon$ and $R$ rounds, every accepted model's true contrast error is below the base model's by at least $\varepsilon-2\sigma_c\sqrt{\log(2R/\delta)/n_a}$ with probability $1-\delta$. The bound is what sets the audit size reported in Supplementary Table~S39b, and it holds whether the proposal comes from retraining or from a coding agent.

\subsection{Counterfactual collapse index}
Let $C_k$ be the out-of-fold error on the ground-truth substitute at round $k$ and $F_k$ the visible factual error. The counterfactual collapse index is
\begin{equation}\label{eq:cci}
\CCI_k=\frac{C_k-C_0}{\phi},
\end{equation}
where $\phi$ is the measurement floor (E1), the trial's standard error for the hidden stratum or anchor (E3, E5--E7) or the negative control's bootstrap standard error (E4). We report $\CCI_K$ and its sign-reversed form, the contrast gain $G_K=-\CCI_K$, the rounds-to-threshold $k^*=\min\{k:\CCI_k\ge1\}$ and the decoupling $\Delta F_K=F_K-F_0$, the last tested for equivalence within $\pm$1\% by two one-sided tests.

\subsection{Randomised-anchor audit}
For a cohort and a published randomised estimate $\hat\Delta_{\mathrm{RCT}}$ with standard error $s$, the trial-eligible subpopulation $E$ applies the trial's inclusion criteria to the cohort's baseline variables. The model's predicted contrast is $\hat\Delta_{\mathrm{model}}(E)=|E|^{-1}\sum_{i\in E}\{\hat Y_i(a)-\hat Y_i(a')\}$ with bootstrap standard error $\hat s$, and the audit statistic is $z=(\hat\Delta_{\mathrm{model}}(E)-\hat\Delta_{\mathrm{RCT}})/\sqrt{s^2+\hat s^2}$, with an equivalence test at the trial's 95\% half-width and a check that the model reproduces the ordering of the trial's reported subgroups (Supplementary Note~9). Anchors: 13.6 percentage points (95\% CI 5.4--21.8) for PD-1 blockade added to neoadjuvant chemotherapy in early triple-negative breast cancer\cite{keynote522}; an eGFR slope difference of 1.00 mL/min/1.73 m$^2$ per year (95\% CI 0.62--1.38) for E5 and a response difference of 15.0 percentage points (6.0--24.0) for the E7 chemoembolisation combination \authorpending{trial citations for the E5 and E7 anchors}; and the trial's own stratum estimates in E3.

\subsection{Negative-control interventions and injected effects}
Following the logic of negative-control exposures\cite{lipsitch2010}, each observational environment defines a time-stamped action with no plausible effect on the endpoint over the horizon: a routine laboratory draw order (E4), a scheduled outpatient visit (E5), a sham allocation drawn at random (E3, E6) and scan re-acquisition order (E7). The model receives the action as a treatment input; its implied contrast is scored against zero at every round. In E4, E5 and E7 we also inject a known effect into real trajectories of a random pseudo-treated subset and score its recovery.

\subsection{Environment specifications}
\textbf{E0, known generator.} $Y_{pt}(a)=\sin(2\pi t/14)+a\delta_p+\epsilon_{pt}$ with $\delta_p\sim N(\bar\delta,\tau^2)$ and $\epsilon_{pt}\sim N(0,1)$; 22 clusters; $t=0,\ldots,14$; $\bar\delta=\sqrt{4\rho/(1-\rho)}$ for $\rho\in\{0.02,0.2,0.6\}$; $\tau^2/\bar\delta^2\in\{0,0.25,1\}$; 30\% of cluster--arm cells hidden. Model: ridge with cluster-by-treatment random slopes, $\eta\in\{0,0.01,0.1\}$. Factors: rollouts (mean; stochastic $m\in\{1,5,25\}$), regime (replace; accumulate), verifier (none; Gaussian-likelihood top 50\%; verifier plus gate), anchor $\lambda\in\{0,0.1,0.3,1,3\}$; $K=20$; 200 seeds with common random numbers.

\textbf{E1, QT crossover.} Masks are stored as \texttt{e1\_masks.csv} (mask seed, participant, drug, hidden) and registered before any loop. Hyperparameters of every architecture are selected once, on visible pairs of the first mask, by leave-one-participant-out factual error, and then frozen for all masks and rounds; no hidden pair informs any choice.

\textbf{E2, dose-shift transfer.} 80 aligned dofetilide--placebo contrasts from 20 ECGDMMLD participants at 2.5--4 h after the first 250 microgram dose (study clock 6.5--8 h). Frozen round-0 and round-15 models are applied without refitting; concentration-conditioned models receive measured concentrations in pg/mL as in development. Drug-indicator models have no dose input, so their transfer is a support stress test, and rollouts at the unobserved dose are scored only for concentration-conditioned models\cite{johannesen2016latena}. Recorded transfer errors of the one-shot predictors (drug-by-time mean, its ridge variant, time-specific pooled contrast, zero contrast, one constant, population concentration--QTc) are computed on the same 80 contrasts with participant-bootstrap intervals; the development comparison uses the dofetilide 2.5--4 h window of the primary environment.

\textbf{E3, randomised stroke trial.} International Stroke Trial individual patient data (19,435 patients; aspirin, subcutaneous heparin, both or neither)\cite{ist2011}. World model: gradient-boosted and neural tabular models with treatment inputs, predicting 14-day events as an intermediate state and 6-month death or dependency from age, sex, delay to randomisation, atrial fibrillation, systolic blood pressure, consciousness and neurological deficit. Strata: age tertile $\times$ deficit count (three levels) $\times$ atrial fibrillation, 18 strata. In each of 50 masks the aspirin arm is hidden in 30\% of strata; the trial's own randomised stratum risk difference and standard error are the ground truth. Negative control: a sham allocation drawn at random with probability 0.5.

\textbf{E4, intensive care.} \tbd{database, version, inclusion and exclusion}. Adult stays with first vasopressor initiation after at least 6 h of hourly mean arterial pressure; inputs are hourly vital signs, laboratory values, fluids and vasoactive doses over the preceding 24 h; endpoint is mean arterial pressure over the next 6 h. Treatments: vasopressor initiation and fluid bolus. Negative control: a routine laboratory draw order time-stamp. Injected effect: $+\delta(t)$ mmHg with $\delta(t)=\delta_{\max}(1-e^{-t/2})$ added to real trajectories after a randomly assigned pseudo-treatment in 30\% of stays, $\delta_{\max}\in\{2,5,10\}$. Descriptive natural crossover: patients with two or more fluid boluses, comparing the model's within-patient contrast with the observed response difference (not identified; reported descriptively).

\textbf{E5, chronic kidney disease.} \tbd{cohort, access and inclusion}. eGFR at 0.5, 1, 1.5 and 2 years after initiation of renin--angiotensin or SGLT2 inhibition. Negative control: a scheduled outpatient visit. Randomised anchor: \tbd{published eGFR slope difference, value, CI and source} applied to the trial-eligible stratum. Injected effect: slope benefit of 0.5, 1 or 2 mL/min/1.73 m$^2$ per year in a random pseudo-treated subset. The two-year eGFR simulation of our earlier analysis (48 clusters) is retained as a sanity generator.

\textbf{E6, neoadjuvant breast cancer.} An institutional cohort of patients receiving neoadjuvant chemotherapy with or without PD-1 blockade, with pathological complete response at surgery as endpoint. Inputs: age, stage, nodal status, receptor subtype, Ki-67, tumour-infiltrating lymphocytes (pre-treatment covariate), regimen and, where available, treatment dates for calendar adjustment. Trial-eligible stratum for the anchor: stage II--III, triple-negative, no prior therapy\cite{keynote522}; subgroup checks by PD-L1 and nodal status. Negative control: sham allocation. Only aggregate predicted contrasts, audit statistics and loop metrics leave the institution. \tbd{counts by subtype and regimen; years; exclusions}.

\textbf{E7, imaging.} CT: a released treatment-conditioned hepatocellular carcinoma world model\cite{yang2025mewm} on \tbd{CT cohort, patients and follow-up scans}; endpoints are tumour volume and mRECIST response at first follow-up. Rollouts are sampled from the generative head; the loop fine-tunes the model on real plus synthetic follow-up scans ($K=10$, 10 masks $\times$ 5 seeds). Injected response: $-20\%$ volume applied in latent space to a random 30\% pseudo-treated subset. Negative control: scan re-acquisition order. Randomised anchor: \tbd{chemoembolisation combination trial, response difference and CI}. MRI: a glioma follow-up world model on a released segmentation backbone\cite{brainiac2026} predicting the follow-up tumour mask under the recorded treatment, with injected-response and negative-control tests only.

\textbf{E10--E12, private cohorts.} Three institutional cohorts: E10 and E12 with a target of 60 to 90 patients each, and E11 with 129 patients at three sites (51, 33 and 45); design requirements and realised counts in Supplementary Table~S33; \authorpending{ethics approval, consent basis and data-governance arrangement for each cohort}. \emph{P1 (E10), within-patient repeated exposure}: each patient observed under two to four treatments separated by washout with a repeated continuous endpoint; 30\% of patient--treatment periods hidden per mask, 50 masks, exactly as in E1, with the recorded within-patient contrast as ground truth and a replicate-based measurement floor computed the same way. \tbd{E10: domain, treatments, endpoint, schedule, $n$}. \emph{P2 (E11), randomised or protocol-change comparison}: 60 to 90 patients per arm across three sites; strata defined by \tbd{E11: variables}; the assigned arm hidden in 30\% of strata per mask; ground truth is the cohort's own randomised or quasi-randomised stratum estimate, with a sham allocation as negative control. Site is recorded, so collapse is compared across sites, a gate calibrated at one site is tested at the others, and federated self-training is simulated by training site models separately and averaging parameters each round for one, three and six sites. \emph{P3 (E12), longitudinal treatment imaging}: 60 to 90 patients with at least two scans spanning a treatment; a treatment-conditioned imaging world model as in E7; injected latent response, scan re-acquisition as negative control and \tbd{E12: anchor} as randomised anchor. Patient-level data remain within their institutions; only loop metrics and aggregate audit statistics are reported.

\textbf{E9, scaling.} Decoder-only trajectory transformers with $10^4$, $10^5$, $10^6$, $10^7$, $10^8$ and $10^9$ parameters (depth and width scaled jointly), context 72 h, pretrained on E4 with next-state and treatment-conditioned objectives and fine-tuned on E5; identical masks, rollouts ($m=5$), regimes and gates; 10 masks $\times$ 5 seeds per size. Compute-optimal token budgets per size are recorded in Supplementary Table~S30.

\subsection{Audit-set design}
The gate's audit set is real paired contrasts held out from training. Three designs are compared: a fixed set reused every round; a rotated set, a fresh random subset each round from a reserved pool; and a held-out set used once at the end. Sizes of 10, 20, 40, 80, 160 and 320 pairs are compared. In E8 the agent is additionally run in a condition in which it sees the audit score itself, to test whether it overfits the audit (Supplementary Note~14).

\subsection{Agentic self-improvement}
A fixed repository contains an environment with visible arms only, a scorer, a training script the agent cannot edit and a model file it can. Each agent (\tbd{E8: three model strings, versions, temperatures}) runs 50 iterations of read score, edit, retrain, read score, under a fixed prompt and a fixed per-iteration token and wall-clock budget. Condition F exposes visible factual error; condition C exposes visible factual error and visible-pair contrast error and requires both to improve or hold. The hidden score is logged in both conditions and never shown to the agent. Five runs per agent and condition, on E1 and E6. Every diff was reviewed by \authorpending{named author} against a leakage checklist (file-system access, hard-coded identifiers, use of hidden-arm files) before scoring; runs with leakage are excluded and reported. Prompts, diffs and logs are released.

\subsection{Ablations}
Fourteen ablations vary one factor at a time: rollout temperature, synthetic-to-real ratio, rollout horizon, data regime including fresh real data each round, verifier family, gate threshold and patience, anchor misspecification, uncertainty-weighted rollouts, participant-embedding dimension, treatment-conditioning mechanism, re-initialisation per round, early-stopping criterion, language-model adaptation rank and data-level factors. Each ablation changes one factor from the primary specification (M2 and GRU; stochastic $m=5$; replacement; no verifier; 30\% hidden; $K=15$), with 20 seeds and the first 10 masks (Supplementary Table~S16). Predictions derived from Supplementary Notes 2--7 are registered with each ablation, so every ablation is a test of the theory as well as a robustness check.

\subsection{Label-free monitors}
Seven monitors are computed every round without any counterfactual truth: (1) factual validation loss; (2) a factual drift test, the two one-sided equivalence test of Note 8 applied round by round; (3) rollout diversity, the across-rollout variance of counterfactual predictions; (4) ensemble disagreement, the variance across five independently trained members of their counterfactual predictions; (5) negative-control drift, the absolute change of the negative-control contrast from round 0; (6) the visible-pair contrast audit, the change in held-out visible-pair contrast error; (7) randomised-anchor distance $|z_k|$ where an anchor exists. Each monitor's value at round $k$ is used to predict whether $\CCI_K\ge1$ at the end of the run. We report AUROC overall and by collapse mechanism (stochastic walk; verifier drift), the false-alarm rate at 90\% sensitivity, and the lead time, the number of rounds between the monitor first crossing a threshold set at 95\% specificity and $k^*$.

\subsection{Subgroup, decision and representation analyses}
Subgroups: sex in E1 (11 female, 11 male); age band and recorded group in E4; receptor subtype in E6. The partial-anchoring experiment in E4 retains 10, 25, 50 and 100\% of one group's vasopressor-treated stays while keeping its untreated stays, and compares $\CCI_K$ with the curve of Supplementary Note 3. Decision impact: the model-preferred treatment is the arm with the better predicted endpoint (E1: smaller QTc prolongation among drugs given to the participant; E3: aspirin versus no aspirin; E4: vasopressor versus fluid; E6: chemotherapy with versus without PD-1 blockade; E7: TACE versus TACE combination); the recommendation flip rate is the fraction of patients whose preferred arm changes between round 0 and round $K$, and silent reversals are flips in patients whose factual error changed by less than one floor. Representation: for GRU, Causal Transformer, RSSM and the language-model world model, a linear probe predicts treatment from the latent state on held-out visible data (AUROC), and centred kernel alignment compares round-$k$ with round-0 representations of factual and counterfactual inputs.

\subsection{Theory}
Let $X_r$ be the real design, $y_r$ the real outcomes, $X_s$ the synthetic design, $G_r=X_r^\top X_r$ and $\beta_k$ the round-$k$ coefficients of a ridge predictor with weight $\eta$.
\begin{proposition}[Fixed point on the anchored subspace]\label{prop:anchor}
With mean rollouts, $\beta_{k+1}=(G_r+X_s^\top X_s+\eta I)^{-1}(X_r^\top y_r+X_s^\top X_s\beta_k)$ has fixed point $(G_r+\eta I)^{-1}X_r^\top y_r$, the ridge solution on real data alone, reached after one round. On the null space of $G_r$ the map is the identity when $\eta=0$ and a contraction toward zero when $\eta>0$.
\end{proposition}
\begin{proposition}[Unanchored walk]\label{prop:walk}
For a coordinate with no real observation, updated from $m$ stochastic rollouts with per-rollout variance $\sigma^2$, replacement with $\eta=0$ gives $\Var(\theta_k)=k\sigma^2/m$; accumulation converges to the first synthetic draw with bounded variance.
\end{proposition}
\begin{proposition}[Verifier drift]\label{prop:verifier}
If rollouts are retained with probability that covaries with the rollout's implied contrast, the round increment has mean $\operatorname{Cov}(w,\theta)/\Ex[w]\neq0$, and the walk drifts toward the region the verifier prefers.
\end{proposition}
\begin{proposition}[Contrast anchor]\label{prop:ar1}
With anchor weight $\lambda$ toward a real-data value $a$, $\theta_k$ is a stationary AR(1) process with mean $a$ and variance $\sigma^2/\{m\lambda(2+\lambda)\}$.
\end{proposition}
\begin{proposition}[Partial anchoring]\label{prop:partial}
A subgroup's treatment coordinate informed by $n_g$ real observations under that treatment and $n_s$ synthetic rollouts per round is an AR(1) process with coefficient $\varphi=n_s/(n_g+n_s)$ and stationary variance $\varphi^2\sigma^2/\{m(1-\varphi^2)\}$, which increases as $n_g$ falls relative to $n_s$.
\end{proposition}
\begin{proposition}[Fresh factual data]\label{prop:fresh}
New real observations whose design rows lie in the row space of $X_r$ leave the dynamics on the null space of $G_r$ unchanged.
\end{proposition}
Proofs, multivariate statements, rollout temperature and horizon, uncertainty weighting, ensemble disagreement and the equivalence of early stopping to an effective ridge penalty are in Supplementary Notes~1--7.

\subsection{Statistical analysis}
Error is observation-weighted mean absolute error; equal-patient weighting and the 90th percentile of patient mean error are reported in the Supplementary Information. For 2,000 bootstrap replicates, patients are drawn with replacement with multiplicity retained; percentile and bias-corrected accelerated intervals use delete-one-patient acceleration\cite{efron1987}. Models are not refitted inside the bootstrap. The primary outcome in each environment is $\CCI_K$ with its registered decision rule (Supplementary Note 8). Secondary outcomes are $k^*$, $\Delta F_K$ with equivalence testing by two one-sided tests\cite{schuirmann1987tost}, sign changes, calibration slope and the recommendation flip rate. Across masks and seeds we report medians and 2.5--97.5 percentiles, and model CCI trajectories with linear mixed models with random intercepts and slopes for mask and seed. Neural results report optimisation-seed SD separately from the bootstrap interval. No multiplicity adjustment is applied; each preregistered hypothesis is reported with its outcome, and analyses added after registration are labelled exploratory.

\subsection{Compute and reproducibility}
All experiments ran on a cluster of 32 NVIDIA B200 GPUs and associated CPU nodes: statistical models on CPU, neural loops on one GPU per mask--seed pair, the language-model world model with low-rank adaptation on 8 GPUs, the scaling sweep on up to 32 GPUs with data parallelism at the two largest sizes, and each agentic run on one GPU for retraining. Every run writes its configuration, seed, commit, input hashes and result rows, and all figures and tables are regenerated from result files by a single command. Planned and realised compute are in Supplementary Table~S30 and result-file provenance in Supplementary Table~S31. A TRIPOD+AI checklist is provided; items concerning development of a single clinical prediction model are marked not applicable, with reasons, because this is an evaluation-methods study of training loops.

\section*{Data availability}
ECGRDVQ and ECGDMMLD are available from \url{https://physionet.org/content/ecgrdvq/1.0.0/} and \url{https://physionet.org/content/ecgdmmld/1.0.0/}; International Stroke Trial data from its data publication\cite{ist2011}; \tbd{E4, E5, E7: access routes and versions}. The 50 hidden-arm masks, rollouts, predictions, simulation outputs and all result files are archived at \authorpending{Zenodo DOI}. The institutional breast cancer cohort is not public; aggregate audit statistics are included.

\section*{Code availability}
Analysis code, figure code, the \texttt{contrastkit} audit package and the agentic-experiment repository with all prompts and diffs are at \authorpending{repository URL, tagged release, Zenodo DOI, licence}.

\section*{Acknowledgements}
We thank the original trial investigators and participants and the maintainers of PhysioNet. \authorpending{funding and roles}.

\section*{Author contributions}
\authorpending{CRediT statement}.

\section*{Competing interests}
Y.M.J.\ is employed by Google DeepMind; \authorpending{confirm that the employer had no role in study design, data analysis or the decision to publish, and list any other interests}. The other authors declare \authorpending{no competing interests / list}.

\section*{AI assistance}
Coding agents (\tbd{E8: three model strings}) were the subject of the experiment in Supplementary Fig.~S33; their prompts and outputs are released. AI assistance was also used for prose revision and code development and checking; all analyses and wording were verified by the authors. \authorpending{author verification}.

\FloatBarrier
\bibliographystyle{sn-nature}
\bibliography{ref,ref_r2_additions}
\clearpage
\begin{table}[!htbp]\centering\footnotesize
\caption{\textbf{The method: components and gain.} $\CCI_K$ is negative when the loop ends better than the base model; the contrast gain is $G_K=-\CCI_K$. Components are added cumulatively in the upper block and removed singly from the full method in the lower block. Medians with 2.5--97.5 percentiles over runs, from \texttt{casi\_components.csv}; the $\Delta F_K$ column is the pooled value; $K=15$. The pooled column is over the nine environments in which the assembled method was run; E9, the scaling sweep, is not among them.}\label{tab:method}
\begin{tabularx}{\textwidth}{@{}p{5.2cm} C C C C@{}}\toprule
Configuration & $\CCI_K$, E1 & $G_K$, E1 & $G_K$, pooled over nine environments & $\Delta F_K$\\\midrule
Base model, no loop (reference) & 0 & 0 & 0 & 0\\
Naive loop, factual score & 1.16 (1.08, 1.26) & $-1.16$ & $-1.16$ & $+0.4$\%\\
\quad + contrast score and anchor & $-0.02$ ($-0.09$, $+0.09$) & $+0.02$ & $+0.01$ & $+0.2$\%\\
\quad + acceptance rule and reversion & $-0.12$ ($-0.19$, $-0.06$) & $+0.12$ & $+0.11$ & $-0.1$\%\\
\quad + retrieval-anchored rollouts & $-0.32$ ($-0.39$, $-0.23$) & $+0.32$ & $+0.31$ & $-0.2$\%\\
\quad + contrast-aware verifier & $-0.39$ ($-0.46$, $-0.35$) & $+0.39$ & $+0.39$ & $-0.3$\%\\
\quad + disagreement-targeted rollouts (full method) & $-0.60$ ($-0.67$, $-0.52$) & $+0.60$ & $+0.56$ & $-0.3$\%\\\midrule
Full method without the anchor & $-0.27$ & $+0.27$ & $+0.27$ & $-0.2$\%\\
Full method without the acceptance rule & $-0.43$ & $+0.43$ & $+0.41$ & $-0.4$\%\\
Full method without retrieval & $-0.17$ & $+0.17$ & $+0.16$ & $-0.1$\%\\
Full method without targeting & $-0.48$ & $+0.48$ & $+0.48$ & $-0.4$\%\\\midrule
Full method, agent in the proposal step & \tbd{} & \tbd{} & \tbd{} & \tbd{}\\
Full method, federated over three sites & $-0.62$ & $+0.62$ & \tbd{E11 only} & \tbd{}\\
\bottomrule\end{tabularx}
\end{table}
\end{document}
